% Generated draft: requires mathematical and source review.
\documentclass[11pt]{article}
\usepackage[margin=25mm]{geometry}
\usepackage{fontspec}
\setmainfont{Latin Modern Roman}
\usepackage{amsmath,amssymb,mathtools,unicode-math}
\setmathfont{Latin Modern Math}
\usepackage{enumitem,needspace,xurl,hyperref}
\setlength{\parindent}{0pt}
\setlength{\parskip}{5pt}
\setlength{\emergencystretch}{3em}
\newcommand{\N}{\mathbb{N}}
\newcommand{\Z}{\mathbb{Z}}
\newcommand{\Q}{\mathbb{Q}}
\newcommand{\R}{\mathbb{R}}
\newcommand{\C}{\mathbb{C}}
\newcommand{\F}{\mathbb{F}}
\newcommand{\A}{\mathbb{A}}
\newcommand{\PP}{\mathbb{P}}
\newcommand{\E}{\mathbb{E}}
\newcommand{\eps}{\varepsilon}
\newcommand{\dd}{\,\mathrm{d}}
\begin{document}
\section*{UnsolvedMath ID: 30001284}
\textbf{Source problem code:} catalogue entry
\subsection{Definitions and mathematical statement}
For a specified base scheme \(S\), write \(\A^1_S=\operatorname{Spec}_S(\mathcal O_S[t])\) for the affine line over \(S\). An \(S\)-group scheme is a group object in the category of \(S\)-schemes. These definitions introduce notation only: the source does not specify a base scheme or its hypotheses, or the intended class of linear algebraic group schemes.

The supplied record contains the following descriptive assertion:
\begin{quote}
For linear algebraic group schemes, the naive algebraic analogue of the classical singular-complex construction computes the \(\A^1\)-homotopy type.
\end{quote}
Here the singular-complex construction, the underlying site and topology, the comparison map, and the meaning of ``computes'' are not specified. In particular, the record does not determine whether additional localization or replacement is required, or which hypotheses on the base and group schemes are intended. Consequently, this assertion cannot be promoted from the supplied evidence to a precise theorem with fixed hypotheses and quantifiers.

\textbf{Unresolved target.} The supplied statement-verification field describes the passage as an assertion containing no question, conjecture, or unresolved task, and marks the statement as unrecoverable. No mathematical open problem can therefore be formulated faithfully from this record alone. The supplied open label and reference to associated comparison questions do not specify such questions.
\subsection{Short English statement}
The record preserves an observation about an algebraic singular construction and \(\A^1\)-homotopy types of linear algebraic group schemes. It supplies neither the hypotheses needed for a precise theorem nor an explicit open question. Its open status is supplied curation, and the intended mathematical target requires human review.
\subsection{Sources}
\begin{itemize}
\item Supplied local registry record: permanent UnsolvedMath ID \(30001284\), title \emph{Descriptive Singular-Complex Observation}, exact source problem code \texttt{catalogue entry}, in the Oberwolfach Reports Open Problems collection. The record lists proposed year 2009 and supplies no named proposer. Its status is \texttt{open}; its statement status is \texttt{unrecoverable}.
\item Supplied report citation: \emph{Algebraic \(K\)-Theory and Motivic Cohomology} (2009), Oberwolfach Reports, pp.\ 1731--1774, DOI 10.4171/owr/2009/31. Source URL and supplied literature-source link labeled ``OWR report'': \url{https://doi.org/10.4171/owr/2009/31}. The supplied page range covers the report; no passage-specific page or individual attribution is given.
\item Supplied literature assessment, checked on \texttt{2026-08-21}: the extracted passage is descriptive and contains no precise resolved problem; associated comparison questions are therefore retained as open. This is supplied dated curation, not an independently verified literature assessment. The background explicitly requests independent verification of time-sensitive status and cited claims. The separate supplied statement-verification field reports that the passage states no question, conjecture, or unresolved task and that no separate clean statement is promoted.
\end{itemize}
\end{document}
